\chapter{Case Study}
\label{chap:experiments}

\section{Model Description}

For the case study, a Probabilistic Risk Assessment (PRA) model for the Generic Modular High Temperature Gas Cooled Reactor (MHTGR) is used \cite{hamza_openpra-orggeneric-mhtgr-model_2025}, which was developed using SAPHIRE and CAFTA. The MHTGR plant comprises four identical 350 MW(t) reactor modules driving two shared turbine generators for a total output of approximately 558 MW(e).

The PRA model encompasses a comprehensive set of event trees and fault trees designed to capture the plant's response to initiating events. The analysis develops seven event trees, each based on a primary initiating event. Table~\ref{tab:mhtgr-initiating-events} shows the corresponding frequencies of the initiating events.

\begin{table}[]
\centering
\caption{MHTGR initiating events and frequencies}
\label{tab:mhtgr-initiating-events}
\begin{tabularx}{\textwidth}{l >{\RaggedRight}X r}
\toprule
\textbf{Initiating Event} & \textbf{Description} & \textbf{Frequency (per year)} \\
\midrule
PCL & Primary coolant leaks & $2.6 \times 10^{-1}$ \\
LOHTL & Loss of main loop cooling & $2.6 \times 10^{0}$ \\
LOOP & Loss of offsite power & $5.0 \times 10^{-3}$ \\
ATRS & Anticipated transients requiring SCRAM & $2.5 \times 10^{1}$ \\
CRW & Control rod withdrawal & $1.0 \times 10^{-1}$ \\
SGTL-S & Small steam generator leaks & $4.0 \times 10^{-1}$ \\
SGTL-M & Moderate steam generator leaks & $4.0 \times 10^{-2}$ \\
\bottomrule
\end{tabularx}
\end{table}

The event tree analysis models accident progression through safety functions including reactor shutdown, decay heat removal, and containment isolation. The complete model comprises 241 event sequences distributed across these 7 event trees as shown in Table~\ref{tab:mhtgr-event-tree-sequences}.

\begin{table}[]
\centering
\caption{Event tree sequences in the MHTGR PRA model}
\label{tab:mhtgr-event-tree-sequences}
\begin{tabularx}{\textwidth}{l >{\RaggedRight}X r}
\toprule
\textbf{Initiating Event} & \textbf{Key Safety Functions} & \textbf{Sequences} \\
\midrule
Small SG tube leaks (SGTL-S) & Moisture detection, SG isolation, decay heat removal & 1 to 46 (46) \\
Primary coolant leaks (PCL) & Reactor trip, RSCE shutdown, SCS/RCCS cooling & 47 to 96 (50) \\
Anticipated transients (ATRS) & SCRAM logic, backup shutdown, heat removal & 97 to 130 (34) \\
Control rod withdrawal (CRW) & PPIS detection, reactor trip, power excursion control & 131 to 138 (8) \\
Loss of offsite power (LOOP) & Reactor trip, backup power, passive cooling & 139 to 150 (12) \\
Moderate SG tube leaks (SGTL-M) & Rapid isolation, pressure control, backup cooling & 151 to 213 (63) \\
Loss of HTS cooling (LOHTL) & Forced convection loss, shutdown, emergency cooling & 214 to 241 (28) \\
\midrule
\textbf{Total} & & \textbf{241} \\
\bottomrule
\end{tabularx}
\end{table}

The fault tree analysis provides detailed reliability models for critical safety systems, with each tree typically decomposed into 5 to 15 subsystems. The primary systems modeled are listed in Table~\ref{tab:fault-tree-systems}.

\begin{table}[]
\centering
\caption{Fault tree models for MHTGR safety systems}
\label{tab:fault-tree-systems}
\begin{tabularx}{\textwidth}{l >{\RaggedRight\arraybackslash}X}
\toprule
\textbf{System} & \textbf{Safety Function} \\
\midrule
Heat transport system (HTS) & Primary heat removal via forced helium circulation \\
Shutdown cooling system (SCS) & Active decay heat removal using auxiliary blower loop \\
Reactor cavity cooling system (RCCS) & Passive heat removal from reactor vessel \\
Reserve shutdown control (RSCE) & Backup neutron absorbing material injection \\
Steam isolation and dump system & Secondary system isolation and pressure control \\
Relief valves & Primary and secondary system pressure relief \\
Reactor trip system & Control rod scram and reactor shutdown \\
%Failure to trip reactor using RSCE (FT34)
%Failure of HTS in at least One Module (FT104)
%Failure of SCS When One Module Requires Cooling (FT174)
%Failure of SCS (FT196)
%Failure of SCS when Four Modules Require Cooling (FT197, FT198 and FT199)
%Intentional depressurization fails (FT215)
\bottomrule
\end{tabularx}
\end{table}

The model includes 170 basic events representing distinct component failure modes and approximately 100 repair events to account for system recovery processes. These events are parameterized with mean failure frequencies or probabilities based on component type and operational context.

\section{Setup Description}
The study used four solver configurations: SAPHSOLVE (Windows version), serial SCRAM, distributed SCRAM, and PRAXIS (in both CPU and GPU modes). All solver configurations were set to compute event sequence frequencies and the number of minimal cut sets for each event sequence. Three truncation settings were evaluated: no truncation, to determine whether all cut sets could be captured, and truncation limits of $10^{-12}$ and $10^{-20}$. Single machine CPU based executions were performed on a system with an Intel i7-13700K processor and 32~GB of RAM under standard processing conditions. Single machine based PRAXIS GPU runs were executed on a workstation equipped with an NVIDIA RTX PRO 6000 Blackwell Max-Q Workstation Edition GPU with 96~GB of GDDR7 memory. Distributed quantifications were executed on a high performance on premise cluster with 12 nodes. For this study, only 5 nodes were used: 1 manager node and 4 worker nodes. The manager and worker nodes are equipped with Intel Xeon CPUs (24 and 128 vCPUs, respectively), 128~GB and 256~GB of RAM, and 4~TB and 1~TB of NVMe storage, respectively. The deployment was scaled to support up to 128 workers.

\subsection{Description of the Studies}
Four studies were performed to evaluate the proposed distributed quantification system. Together, they demonstrate the capabilities and advantages of the system, highlight its limitations, and examine the effectiveness of the measures introduced to mitigate those limitations.

(a) \textbf{Multiple concurrent job handling:}
The first study evaluates how well the distributed system handles large numbers of concurrent jobs submitted by users. It focuses on scalability and speedup as the number of workers increases. The CRW event tree is quantified using three configurations: SAPHSOLVE (serial), SCRAM (serial), and distributed SCRAM. For distributed SCRAM, the number of workers is varied from $2$ to $128$, while the number of simultaneously submitted jobs is varied from $1$ to $10,000$.

(b) \textbf{Robust execution of large event tree jobs:}
Very large quantification tasks can run for a long time, become stuck, or fail after many hours, making it difficult to determine whether they will ever complete and risking complete loss of computational effort upon failure. The second study addresses this by introducing a three stage workflow (preprocessing, processing, and postprocessing) in which a large event tree is decomposed into individual event sequences. This decomposition enables partial retention of results, even if some sequences fail, and allows the distributed system to requeue failed jobs or to move repeatedly failing jobs into dead letter queues, and to continue processing the remaining sequences. In this way, the system aims to guarantee execution of large event tree jobs in a fault tolerant manner, or at least to preserve partial results under worker or job failures. For this study, the MHTGR model is quantified using SAPHSOLVE, serial and distributed SCRAM with $2$ to $128$ workers.

(c) \textbf{Stress testing:}
The final study focuses on the overall throughput and latency of the distributed system under large request loads. The LOOP event tree is used with a truncation limit of $10^{-12}$. This model is repeatedly submitted to the quantification API endpoint up to $10^{6}$ times, and an equal number of requests is then issued to retrieve the corresponding results. The number of workers is fixed at $128$ for this stress test. The objective is to determine how many quantification requests and result-retrieval operations the system can handle concurrently without failure, and to characterize the limits at which throughput and latency begin to degrade, even for relatively small jobs with a large worker pool.

The next section presents the detailed description and the results of these four studies.

\section{Results and Discussion}
\subsection{Concurrent Job Handling and Performance}
This study evaluates how well distributed SCRAM handles large numbers of concurrent jobs and how its performance changes as the workload increases. Practical PRA studies can contain tens of event trees, tens to hundreds of fault trees, and hundreds of event sequences. Repeated calculations of this kind also appear in sensitivity, importance, and uncertainty analyses, where the same quantification must be performed many times while only the probabilities of basic events or CCF groups are changed.

The CRW event tree from the MHTGR model was used as the benchmark model. One full evaluation of the CRW event tree was counted as one job. Since the CRW event tree contains 8 event sequences, the workload range of 1 to 10{,}000 jobs corresponds to 8 to $8.00E+04$ event sequences in total. A comparison is made between serial SAPHSOLVE (Windows version), single-process SCRAM, and distributed SCRAM configurations using 2, 4, 8, 16, 32, 64, and 128 workers. As a representative example, 10{,}000 jobs correspond to changing the probabilities of 1{,}000 basic events from 0.1 to 1.0 in increments of 0.1.

\begin{figure}[htbp]
  \centering
  \includesvg[width=1\textwidth]{6-experimental-evaluation/figures/performance.svg}
  \caption{Execution time of as a function of the number of submitted jobs (from $1$ to $1000$) for SAPHSOLVE, SCRAM, and distributed SCRAM with 2 to 128 workers.}
  \label{fig:performance}
\end{figure}
\begin{figure}[htbp]
  \centering
  \includesvg[width=1\textwidth]{6-experimental-evaluation/figures/speedup.svg}
  \caption{Speedup relative to SCRAM Serial as a function of the number of concurrent jobs (from $1$ to $10000$) for SAPHSOLVE, SCRAM, and distributed SCRAM with 2 to 128 workers.}
  \label{fig:speedup}
\end{figure}

Figure~\ref{fig:multiple_concurrent_job} shows the wall-clock time as a function of the number of jobs and workers, and figure~\ref{fig:performance} shows the speedup of the distributed system relative to the serial solvers for different numbers of workers. At the low end of the workload range, the fixed cost of job creation, queueing, broker communication, dispatch, worker coordination, and result collection dominates. For this reason, when only a few jobs are submitted, single process SCRAM is usually faster than the distributed configurations, and increasing the number of workers beyond the number of available jobs provides little benefit. As the number of jobs increases, this fixed overhead is amortized, the workers remain busy for longer periods, and the benefit of distributed execution becomes clear.

For $1{,}000$ jobs, SAPHSOLVE required about $9.31E+03$~seconds ($\approx 1.55E+02$~minutes), while single process SCRAM required about $6.57E+03$~seconds. Distributed SCRAM reduced the runtime to about $3.53E+03$, $1.99E+03$, $8.53E+02$, $4.29E+02$, $2.16E+02$, $1.15E+02$, and $6.60E+01$~seconds for 2, 4, 8, 16, 32, 64, and 128 workers, respectively. For $10{,}000$ jobs, SAPHSOLVE required about $9.31E+04$~seconds ($\approx 2.59E+01$~hours), while single process SCRAM required $6.56E+04$~seconds ($\approx 1.82E+01$~hours).

The speedup $S(N)$ for $N$ workers is defined in the standard way as:
\begin{equation}
  S(N) = \frac{T_{\text{serial}}}{T_{\text{parallel}}(N)}
\end{equation}
where $T_{\text{serial}}$ is the wall-clock time of the single process SCRAM run, and $T_{\text{parallel}}(N)$ is the wall-clock time when the same workload is executed using $N$ workers in the distributed system. Using distributed SCRAM, the wall-clock time was reduced to approximately $3.43E+04$, $1.70E+04$, $8.53E+03$, $4.29E+03$, $2.17E+03$, $1.14E+03$, and $6.45E+02$~seconds for 2, 4, 8, 16, 32, 64, and 128 workers, respectively. Taking the single process SCRAM runtime as the serial baseline for the $10{,}000$ job case, the corresponding speedups are 1.91x, 3.85x, 7.69x, 15.30x, 30.16x, 57.52x, and 101.67x.

Even though the jobs in this experiment are independent and, in principle, $100\%$ parallelizable, the scaling curves still show behavior similar to Amdahl's law. The computational part of each job is embarrassingly parallel, but the full workflow still contains serial and overhead components that do not disappear when more workers are added. Jobs must be created by the producer, submitted to the message broker, placed in a queue, dispatched to workers, and collected again when the results are returned. Each worker must receive a job, deserialize its input, initialize or load the model if necessary, execute the analysis, serialize the results, and send them back. These additional steps introduce time that cannot be parallelized.

As the number of jobs increases, the system initially shows near linear scaling, but the speedup eventually begins to flatten. For a fixed number of workers, increasing the number of jobs increases queueing delay (latency) because more jobs must wait before a worker becomes available. For a fixed workload, adding more workers reduces the computation time per worker, but it does not remove the above mentioned overheads. As a result, these costs become a larger fraction of the total runtime at high worker counts.

This explains why, even with 128 workers, the maximum observed speedup for $10{,}000$ jobs is 101.67x rather than the ideal 128.00x. The gap can be attributed to the time required for the producer to package and enqueue each job, network latency between the producer, the message broker, and the workers, internal routing delay inside the message broker, the cost of serializing and deserializing job data, worker side model initialization, and contention for shared resources such as disks or file systems when writing logs or saving results.

An important feature of this study is that the resulting cut sets remain the same and only the overall probabilities or frequencies change. Because of this, the runtime of each job is very similar. This symmetric workload is easy to balance across workers and contributes to the good, though not perfectly linear, scaling observed in the figures. Overall, the results show that distributed SCRAM can substantially reduce the execution time of large batches of independent PRA calculations, including sensitivity analyses, importance measure studies, uncertainty quantification, and repeated scenario evaluations, while also showing that communication and coordination overhead place a practical upper bound on the achievable speedup.

\subsection{Robust Execution of Large Event Tree Jobs}
Very large event tree quantification jobs are difficult to execute robustly because a single stalled or failed run can waste many hours of computation and provide little information about which part of the model caused the problem. To study this behavior, the full set of MHTGR event trees was quantified at a truncation limit of $1E-20$ using SAPHSOLVE and SCRAM. In the baseline run, each event tree was submitted as a single large job without event sequence decomposition. Under this setup, both tools were able to solve only 5 of the 7 event trees. In particular, the PCL and ATRS event trees could not be solved. The baseline results have been summarized in Table~\ref{tab:mhtgr_robust_run1}.

\begin{table}[]
\centering
\caption{Baseline event tree level quantification results for the seven MHTGR event trees at a truncation limit of $10^{-20}$.}
\label{tab:mhtgr_robust_run1}
\small
\begin{tabular}{@{}lrrr@{}}
\toprule
\textbf{Event tree} & \textbf{Minimal cut sets} & \textbf{SAPHSOLVE (s)} & \textbf{SCRAM (s)} \\
\midrule
SGTL-S & $1.75 \times 10^{5}$ & $1.37 \times 10^{1}$ & $3.75 \times 10^{1}$ \\
PCL & N/A & N/A & N/A \\
ATRS & N/A & N/A & N/A \\
CRW & $3.25 \times 10^{6}$ & $1.76 \times 10^{2}$ & $9.24 \times 10^{1}$ \\
LOOP & $4.34 \times 10^{4}$ & $4.02 \times 10^{0}$ & $5.81 \times 10^{1}$ \\
SGTL-M & $2.33 \times 10^{5}$ & $2.02 \times 10^{1}$ & $5.42 \times 10^{1}$ \\
LOHTL & $5.36 \times 10^{6}$ & $2.53 \times 10^{2}$ & $4.21 \times 10^{2}$ \\
\bottomrule
\end{tabular}
\end{table}

The same calculation was then repeated in the distributed system using the SAPHSOLVE and SCRAM executables via the \texttt{/q/execute/saphsolve} and \texttt{/q/execute/scram} endpoints, with event sequence decomposition enabled. For SAPHSOLVE, this was done by splitting the input file into separate event sequence inputs, since SAPHSOLVE is closed source and only its file level interface can be manipulated. For SCRAM, the PDAG construction step was decomposed. In the original executable workflow, individual PDAGs are created and quantified in a "for" loop. In the distributed workflow, the PDAGs for all event sequences are obtained first and then distributed to workers, which quantify them independently and return sequence level results for final aggregation. The results of this second run have been summarized in Table~\ref{tab:mhtgr_robust_run2}.

\begin{table}[]
\centering
\caption{Sequence decomposed execution results for the seven MHTGR event trees at a truncation limit of $10^{-20}$.}
\label{tab:mhtgr_robust_run2}
\small
\begin{tabular}{@{}lrrr@{}}
\toprule
\textbf{Event tree} & \textbf{Minimal cut sets} & \textbf{SAPHSOLVE (s)} & \textbf{SCRAM (s)} \\
\midrule
SGTL-S & $1.75 \times 10^{5}$ & $1.55 \times 10^{1}$ & $3.75 \times 10^{1}$ \\
PCL & $8.30 \times 10^{8}$ & $1.51 \times 10^{5}$ & $1.90 \times 10^{4}$ \\
ATRS & $2.10 \times 10^{8}$ & $2.46 \times 10^{4}$ & $1.89 \times 10^{3}$ \\
CRW & $3.25 \times 10^{6}$ & $2.18 \times 10^{2}$ & $9.96 \times 10^{1}$ \\
LOOP & $4.34 \times 10^{4}$ & $4.28 \times 10^{0}$ & $5.83 \times 10^{1}$ \\
SGTL-M & $2.33 \times 10^{5}$ & $2.22 \times 10^{1}$ & $5.43 \times 10^{1}$ \\
LOHTL & $5.36 \times 10^{6}$ & $2.91 \times 10^{2}$ & $4.22 \times 10^{2}$ \\
\bottomrule
\end{tabular}
\end{table}

A direct comparison of the two result summaries shows the main trade-off of this approach. For event trees that were already solvable in the baseline run, sequence level decomposition adds a small amount of overhead because the system must split the work, dispatch more jobs, and aggregate the returned results. However, this overhead is outweighed by the gain in robustness for difficult jobs. In the decomposed workflow, SCRAM was able to quantify all event trees, including PCL and ATRS, which failed in the original run. For SAPHSOLVE, ATRS also became solvable, and the PCL failure no longer appeared as an unexplained full job failure. Instead, the problematic sequence was isolated to PCL-63.

This is the main advantage of the preprocessing, processing, and postprocessing workflow. Failed sequence jobs can be re-queued for quantification automatically, and if a sequence keeps failing it can be moved to a dead letter queue while the remaining event sequences continue to run. As a result, the full run does not lose all completed work, and the user still receives the partial results together with a clear indication of which sequence caused the problem. In this study, PCL-63 continued to fail in SAPHSOLVE and was therefore moved to the dead letter queue, while the remaining sequences completed normally.

In the decomposed run, the two dominant event trees were PCL and ATRS. PCL generated $8.30E+08$ minimal cut sets and required approximately $4.20E+01$~h in SAPHSOLVE and $5.28E+00$~h in SCRAM, while ATRS generated $2.10E+08$ minimal cut sets and required approximately $6.84E+00$~h in SAPHSOLVE and $5.25E-01$~h in SCRAM. At the event sequence level, the three largest contributors were PCL-91, PCL-76, and PCL-63, which generated approximately $3.67E+08$, $2.23E+08$, and $1.99E+08$ minimal cut sets, respectively. This explains why the PCL event tree dominates the total execution time and why sequence level distribution is especially useful for such workloads.

The main points of this calculation are summarized below:
\begin{itemize}
    \item Sequence level decomposition changed the previously unsolved PCL and ATRS event trees into solvable workloads in SCRAM.
    \item For SAPHSOLVE, decomposition exposed the main failure point as PCL-63 instead of losing the entire PCL job without diagnosis.
    \item The distributed workflow improved robustness by allowing successful sequences to finish, failed sequences to be retried, and repeatedly failing sequences to be moved to a dead letter queue.
    \item The approach also revealed the main computational hotspots, namely PCL-91, PCL-76, and PCL-63, which produced the largest numbers of minimal cut sets and dominated the runtime.
\end{itemize}

Overall, these results show that event sequence decomposition is not only a parallelization strategy but also a robustness mechanism. It makes large event tree jobs easier to diagnose, more fault tolerant during execution, and less likely to lose all computational effort when part of the workload fails.

\subsection{Load Testing / Stress Testing}

\begin{figure}[htbp]
  \centering
  \includesvg[width=1\textwidth]{6-experimental-evaluation/figures/load_test_analysis.svg}
  \caption{Performance analysis of microservices deployment showing (a) heavy batch processing scaling, (b) high concurrency limits, and (c) sustained load stability.}
  \label{fig:load_test_analysis}
\end{figure}

\subsubsection{Test 1: Heavy batch processing (LOOP Model)}
In this scenario, a workload consisting of 10,000 users submitting heavy batch jobs ($>0.5MB$), each requiring the decomposition of event trees into individual sequences, was simulated. As shown in \textbf{Figure 1(a)}, scaling the number of job brokers from 1 to 6 resulted in a linear increase in throughput (from $25.6$ to $132.5$ req/s) and a significant reduction in latency (from $2.52$s to $0.48$s). However, increasing the replicas to 8 caused a sharp degradation in performance, with throughput dropping to $38.9$ req/s. This indicates a bottleneck caused by connection instability rather than computation limits. Specifically, the intense traffic saturated the message queueing and object storage services (RabbitMQ and MinIO), preventing them from responding to health check pings; this triggered connection drops after two consecutive health check failures. Consequently, \textbf{6 brokers} was identified as the optimal configuration for such heavy workload.

\subsubsection{Test 2: High concurrency limits (Chinese Model)}
For lightweight, non-batch jobs ($\sim 3KB$), the system limits were evaluated by increasing the number of concurrent users from 10,000 to 25,000. \textbf{Figure 1(b)} illustrates that the system maintained a $0\%$ error rate up to 15,000 users. However, pushing beyond this threshold to 20,000 users introduced a $1.13\%$ error rate, despite the throughput remaining relatively stable. Furthermore, scaling brokers from 8 to 64 at the safe limit of 10,000 users yielded negligible gains, plateauing at $\sim1200$ req/s. This suggests that for lightweight jobs, the bottleneck again shifts downstream to the ingestion capacity of the message broker and object storage services (RabbitMQ/MinIO).

\subsubsection{Test 3: Sustained load stability}
To assess long-term stability, the system was subjected to a sustained load of 10,000 users for 10 minutes. \textbf{Figure 1(c)} demonstrates that the system remained highly stable across 8, 16, and 64 broker configurations, processing over 1.4 million requests without failure. While the 64-broker setup provided the lowest latency ($3.97$s) and highest throughput ($1212$ req/s), the 8-broker configuration was surprisingly competitive, recording $4.18$s latency and $1195$ req/s throughput. This reinforces the conclusion that for sustained, lightweight workloads over-provisioning brokers provides diminishing returns.

\begin{table}[htbp]
    \centering
    \caption{Summary of load testing results and bottleneck analysis}
    \label{tab:load-test-summary}
    \begin{tabular}{@{}llll@{}}
        \toprule
        \textbf{Test ID} & \textbf{Workload Type} & \textbf{Optimal Config} & \textbf{Primary Bottleneck} \\ 
        \midrule
        Test 1 & Heavy Batch ($>0.5$ MB) & 6 Brokers & Brokers and service connection timeouts \\
        Test 2 & High Concurrency ($\sim3$ KB) & 15,000 Users (Max) & Brokers and service connection limits \\
        Test 3 & Sustained Load (10 min) & 8 Brokers & N/A (diminishing returns) \\
        \bottomrule
    \end{tabular}
\end{table}

\subsection{Cost Estimate for Cloud Based Deployment}
\begin{figure}[H]
  \centering
  \includesvg[width=\textwidth]{6-experimental-evaluation/figures/aws.svg}
  \caption{High level overview of the AWS infrastructure deployment for PRA quantification purposes.}
  \label{fig:aws}
\end{figure}

To place the performance and cost profile of the distributed system presented in this work in perspective, the costs associated with running an identical workload on a state-of-the-art commercial cloud platform were analyzed. The projected operational expenditure for such an AWS hosted environment is estimated at approximately \$11,413 per month. The primary cost drivers are security governance (IAM Access Analyzer) and high volume storage (S3), which together account for over 55\% of the monthly budget. This estimate assumes a consistent workload utilizing 128 instances and significant message brokerage capacity to handle the load defined in the previous stress tests.

To derive these figures, usage patterns were modeled based on a user base of 10,000 active participants. The API Gateway costs were calculated assuming each user operates at the rate limit of 10 requests per hour continuously (24 hours over 30 days), resulting in a total of 72 million monthly requests. Identity management (IAM) costs reflect the maximum allowance of 10,000 distinct monitored accounts. Compute and messaging resources were provisioned to match the stress test environment, consisting of a single EKS cluster orchestration plane managing 128 EC2 worker nodes (1 vCPU, 4 GB RAM each) and 8 dedicated Amazon MQ broker instances. Storage costs for S3 are projected based on a hard cap of 10 GB per user, totaling 100 TB of persistent data.

It is important to note that several managed AWS services were explicitly excluded from this estimate as their functionality is handled by platform-agnostic alternatives in our architecture. Specifically, external container registries (Docker Hub/GitHub Container Registry) are used instead of Amazon ECS, dead letter queues (DLQs) are managed natively within RabbitMQ, autoscaling is handled through KEDA, and MongoDB is used in place of Amazon DynamoDB.

\begin{table}[htbp]
    \centering
    \caption{AWS Infrastructure Cost Estimate (US East - Ohio)}
    \label{tab:aws-cost-estimate}
    \begin{tabular}{@{}llrr@{}}
        \toprule
        \textbf{Service} & \textbf{Configuration Details} & \textbf{Monthly (USD)} & \textbf{Annual (USD)} \\ 
        \midrule
        IAM Access Analyzer & 10,000 accounts monitored & 4,000.00 & 48,000.00 \\
        S3 Standard & 100 TB storage, Data transfer & 2,321.74 & 27,860.88 \\
        Amazon EC2 & 128 $\times$ m6g.medium & 1,812.74 & 21,752.83 \\
        Amazon MQ & 8 $\times$ m5.large (RabbitMQ) & 1,784.32 & 21,411.84 \\
        Amazon CloudWatch & Metrics \& DB Insights & 1,168.90 & 14,026.80 \\
        Amazon API Gateway & 72M Requests/mo & 252.00 & 3,024.00 \\
        Amazon EKS & 1 Cluster & 73.00 & 876.00 \\
        \midrule
        \textbf{Total} & & \textbf{\$11,412.70} & \textbf{\$136,952.40} \\
        \bottomrule
    \end{tabular}
\end{table}